\chapter{函数空间.}

我们知道，任意函数的概念大致直到 19 世纪初才出现。至于以一般的方式研究函数集合、并为它们配备拓扑结构的思想，则不早于黎曼（见第 141 页）出现，而且直到 19 世纪末才开始付诸使用，就更是如此了。

然而，数值函数序列收敛的概念，从无穷小演算的开端起就或多或少自觉地被使用着。但那只是简单收敛的问题，而且在波尔查诺和柯西精确地定义收敛级数和连续函数的概念之前，也不可能不是如此。后者最初并没有看出简单收敛与一致收敛之间的区别，还以为能够证明每一收敛的连续函数项级数都以连续函数为其和（{[}56 a{]}，(2)，第 III 卷，第 120 页）。这一错误几乎立即被阿贝尔揭露，他同时证明了每个整幂级数在其收敛区间的内部是连续的，所用的论证已成为经典，它在这一特殊情形本质上用到一致收敛的观念（{[}1{]}，第 I 卷，第 223–224 页）。剩下的只是以一般的方式把这一观念分离出来，这由斯托克斯和塞德尔在 1847–48 年、以及柯西本人在 1853 年（{[}56 a{]}，(1)，第 XII 卷，第 30 页）独立地完成。\(^{1}\)

在魏尔斯特拉斯和黎曼的影响下，一致收敛概念及相关问题的系统研究在 19 世纪最后三分之一的时间里由德国学派（汉克尔、杜布瓦–雷蒙）尤其是意大利学派发展起来：迪尼和阿尔泽拉明确了连续函数序列的极限为连续的必要条件，而阿斯科利则引入了等度连续性这一基本概念，并证明了刻画连续函数之紧集的定理{[}10{]}（这个定理后来由蒙泰尔在他的“正规族”理论中加以普及，正规族无非就是解析函数的相对紧集）。

另一方面，魏尔斯特拉斯本人（{[}329 a{]}，第 III 卷，第 1–37 页）发现了用多项式在一个有界集上一致逼近一个或多个实变量的连续数值函数的可能性，这一结果立即激起了强烈的兴趣，并引出了许多“定量”研究，它们处于我们目前所遵循的观点之外。

这些问题上的现代贡献，首先在于赋予它们以所能有的全部广度，即把它们用于定义集和值集不再局限于 R 或有限维空间的函数，并借助一般的拓扑概念把它们安置到其自然的框架之中。特别地，已在古典分析中显示出是一阶工具的魏尔斯特拉斯定理，近年来被 M. H. 斯通推广到了远为一般的情形：通过发展 H. 勒贝格引入的一个想法（在魏尔斯特拉斯定理的一个证明中），他突显了格在连续数值函数逼近中所起的重要作用（用“格多项式”逼近），另一方面又表明魏尔斯特拉斯定理的推广形式如何立即蕴涵一整串类似的逼近定理，它们就这样以一种连贯得多的方式被归拢在一起{[}301 c{]}。

